% Plan: development/outline.md, section 5, "§11". Round-1 revision
% (adjudication G4-01, G4-08, G4-09; LD-1(b), LD-5): AI-use statement added to
% the declarations; \archiveDOI and \TODO mark the remaining declaration
% placeholders (decision D8, R-12). Limitations point to the single statements
% in Sections 2 and 3 and to thm:eg-bounds (eg trust base).
% Numbers not in data/numbers.json: 0.61 (max of D_n/(n^2 eps) = 0.605 in
% tab:camshape-deficit, rounded up); 87% (camshape register CS-05);
% 37 = 31 closures + 5 waterno2 + ann;
% 114 = 155 - 29 - 12 (tab:results-funnel); 15 instrumented runs and 122
% runs (obs:scip-propagation); 15 staged closures = 11 splits (lnts x4, dtoc5,
% optcdeg2, lukvle10, chain x4) + 4 catmix value-function minorants.
% Round-2 revision (adjudication G4-06 to G4-08; decisions LD2-1, LD2-2): the
% recommendations moved to sec:conclusion-recs in 09-interpretation.tex (move
% M9). 27 of 31 = 12 + 6 + 7 + 2 closures with branching absent or at most
% three-dimensional (sec:interpretation-structure); pindyck and the three eg
% instances are the other 4. Six families with a weaker or partial second
% implementation: optcdeg2, lukvle10, catmix, etamac, pindyck (weaker second)
% and eg (partial second), tab:trust. AI-use statement: same wording as
% sec:semantics-protocol (adjudication G1-01); evidence in adjudication section
% 1 (root-research-log, solver-runs report, data-semantics.md).
% Round-3 revision (development/reviews/round3/), checked 2026-10-05: proofs
% reread by separate agent sessions: Appendix B (round-1 sol-math-main and
% opus-math-main; development/open-items.md item 1), chain and catmix
% (B4-chain-catmix app:chain-prior), an earlier form of the eg rounding
% analysis (F-eg-rounding app:egrounding-trust), eg primal check (C-points;
% B7-eg app:eg-primal). 114 = 155 - 41 (29 closures + 12 other instances of
% ours among the 155 that pass the selection funnel).
\section{Limitations and conclusion}\label{sec:conclusion}

\paragraph{Limitations.}\label{sec:conclusion-limits}
Nothing in this paper is formally verified; \cref{sec:semantics-trust,tab:trust} state what each certificate trusts, and the \inst{eg} dual bounds also assume correct code and the faithful execution and retention of the recorded run (\cref{thm:eg-bounds}).
All implementations come from agent sessions directed by the authors; for six certificate families the second implementation certifies only a slightly weaker dual bound or covers only part of the domain, and for \inst{ann_cumene_tanh} it is not separately written; \cref{sec:semantics-protocol} states these limits and the resulting common-mode risk.
Separate agent sessions reread the proofs of \cref{app:splitproofs}, the \inst{chain} and \inst{catmix} proofs (\cref{app:chain-prior}) and an earlier written form of the \inst{eg} rounding analysis (\cref{app:egrounding-trust}), and reran and read the \inst{eg} primal check (\cref{app:eg-primal}); no person outside the authors has checked the code or the proofs (\cref{sec:semantics-protocol}).
Replay was run on one machine.
The \evid{rerun} certificates store no per-leaf or per-stage certificate data; apart from the \inst{eg} leaf partitions and the saved node counts from which the \inst{ann_cumene_tanh} replay reconstructs its partition, their search partitions are not stored either.

Unless a statement names another reading, the claims concern \reading{b} of the stored OSIL files (\cref{sec:semantics-model} names the exceptions), and statements about listed data concern the pages fetched on the dates of \cref{sec:semantics-model}; \cref{rem:sem-readings} states what transfers to the GAMS text and to the binary64 data that solvers read.
The certificates were built individually, one instance or family at a time; we implemented no general solver capability.
We chose the instances by judged tractability and did not attempt 114 of the 155 instances that pass our selection funnel, so our counts do not measure a solve rate (\cref{sec:results-selection}).
Some sources could not be read \citep{mcdonald1997-glopeq-a-new-computational-tool,huang2011-operative-planning-of-water-supply,ghaddar2015-a-lagrangian-decomposition-approach-for}, the first of them cited by the MINLPLib pages of \inst{ex6_2_5} and \inst{ex6_2_7}; our priority statements are limited accordingly (\cref{app:literature-unread}).

\paragraph{Open problems and conclusion.}\label{sec:conclusion-open}
The five \inst{waterno2} instances (relative gaps of 1.68\% to 10.82\%) and \inst{ann_cumene_tanh} (0.195\%) remain unclosed.
Closing KAN instances whose stored B-spline data satisfy the partition-of-unity rows exactly is open, and so is whether the GLOPEQ results of McDonald and Floudas cover \inst{ex6_2_5} and \inst{ex6_2_7}.
The short certificates of \inst{camshape}, \inst{dtoc5} and \inst{hvycrash} are natural candidates for a check in a proof assistant.

We return to the three observations of \cref{sec:intro}.
First, once a bound fitted to the structure of the model was in place, the remaining branching was low-dimensional for 27 of the 31 closures; we have only interpreted, not tested, why general-purpose solvers did not find such bounds (\cref{sec:interpretation}).
Second, fifteen closures exploit the stage structure, eleven through splits and four through value-function minorants.
Third, proving exact feasibility settled conflicts among listed bounds and changed the verdict on listed points and solver claims often enough that benchmark values should carry their data semantics, their provenance and, where possible, a certificate that others can check.

\section*{Statements and declarations}

\paragraph{Data and code availability.}
Code, certificate data, exact point definitions, logs, the number file and the claim register form the archive that accompanies this paper; it will be deposited at \archiveDOI{} under \TODO{licence of the archived code and data}.
The archive includes MINLPLib and QPLIB models under their CC~BY~4.0 licence \citep{vigerske2026-minlplib-a-library-of-mixed,furini2018-qplib-a-library-of-quadratic}; CUTEst and SIF problem files, where the archive includes them, keep their own notices (three-clause BSD for the ralna/SIF collection, MIT for the optrove/sif copies).
The MATPOWER case data behind the \inst{powerflow} models \citep{zimmerman2011-matpower-steady-state-operations-planning}, papers and other copyrighted sources are not redistributed; the archive lists the SHA-256 hashes of the copyrighted sources.

\paragraph{Use of AI tools.}
AI agents (Anthropic Claude and OpenAI GPT models), working under the authors' direction, selected the instances, devised the certificate constructions, wrote the proofs, implemented the computations and re-implemented them separately, checked code and proofs, ran the solver experiments and the audit, searched the literature and drafted the manuscript.
\Cref{sec:semantics-protocol} states how the implementations were produced and checked.
The authors take full responsibility for the content.
\TODO{authors to confirm, here, in \cref{sec:semantics-protocol} and in the limitations of \cref{sec:conclusion}, including which proofs and code the authors checked themselves}

\paragraph{Competing interests.}
\TODO{competing interests; if none: ``The authors have no competing interests to declare.''}

\paragraph{Funding.}
\TODO{funding statement; if none: ``No funding was received for this work.''}
